\section{Chance-constrained validator-set sizing}
\label{sec:sizing}

\subsection{Problem statement}

\begin{Problem}[Validator-set sizing]
\label{prob:sizing}
Given a peak demand $\Dpk$, an achievable client concurrency $\cc$, a
performance risk level $\epsp$, a safety risk level $\epss$, a detection window
$\w$, and a cost function $\Cost:\NN\to\RR_{+}$, find
\begin{equation}
\nopt\in\arg\min_{\n\in\NN,\ \n\ge4}\ \Cost(\n)
\qquad\text{subject to}\qquad
\begin{cases}
\Prob\bigl(\TPS(\n,\cc)\ge\Dpk\bigr)\ \ge\ 1-\epsp, & \text{(P)}\\[0.6ex]
\pmiss(\n,\w)\ \le\ \epss. & \text{(S)}
\end{cases}
\label{eq:sizing}
\end{equation}
\end{Problem}

Constraint~(P) is a chance constraint in the sense
of~\cite{charnes1959}: the calibrated model is uncertain, and we require the
demand to be met with confidence $1-\epsp$ rather than in expectation.
Constraint~(S) is the detection requirement of
Corollary~\ref{cor:nmin}. The lower bound $\n\ge4$ is the smallest set for
which $\lfloor(\n-1)/3\rfloor\ge1$.

\subsection{Deterministic equivalent}

\begin{Lemma}[Deterministic equivalent of the chance constraint]
\label{lem:detequiv}
Under the predictive distribution of Statement~\ref{st:pi},
constraint~(P) holds if and only if
\begin{equation}
h(\n)\eqdef\log\Tzero+\log\gfun(\cc)-\beta\log\n
-\zq{1-\epsp}\,s(\n,\cc)-\log\Dpk\ \ge\ 0 .
\label{eq:detequiv}
\end{equation}
\end{Lemma}

\begin{Proof}
$\log\TPS(\n,\cc)$ is Gaussian with mean $\hat y=\log\Tzero+\log\gfun(\cc)
-\beta\log\n$ and standard deviation $s(\n,\cc)$. Hence
$\Prob(\TPS\ge\Dpk)=\Prob\bigl(Z\ge(\log\Dpk-\hat y)/s\bigr)
=\Phi\bigl((\hat y-\log\Dpk)/s\bigr)$, and this is at least $1-\epsp$
exactly when $(\hat y-\log\Dpk)/s\ge\zq{1-\epsp}$, which
is~\eqref{eq:detequiv}.
\end{Proof}

\begin{Lemma}[Strict monotonicity]
\label{lem:monotone}
If $\beta>0$ and $s(\n,\cc)$ is nondecreasing in $\n$ on the extrapolation
range, then $h$ is strictly decreasing in~$\n$.
\end{Lemma}

\begin{Proof}
$\partial h/\partial\n=-\beta/\n-\zq{1-\epsp}\,\partial s/\partial\n<0$, since
$\beta>0$, $\n>0$, $\zq{1-\epsp}>0$ for $\epsp<1/2$, and
$\partial s/\partial\n\ge0$.
\end{Proof}

The condition $\partial s/\partial\n\ge0$ holds automatically for the variance
form~\eqref{eq:predvar} whenever $\log\n_{t}$ lies outside the calibration
range, which is the case of interest.

\subsection{The feasibility window}

\begin{Theorem}[Feasibility window]
\label{thm:window}
Let $\beta>0$, $\epsp<1/2$, $\varepsilon_{0}>0$, and let $s(\cdot,\cc)$ be
nondecreasing. Define
\begin{align}
\nmax &\eqdef\sup\{\n>0:\ h(\n)\ge0\},
\label{eq:nmax}\\
\nmin &\eqdef\frac{\bigl(1+(\kf-1)\rho\bigr)\ln(1/\epss)}
                  {2\,\w\,\varepsilon_{0}^{2}} .
\label{eq:nmin-repeat}
\end{align}
Then the feasible set of Problem~\ref{prob:sizing} is the closed integer
interval
\begin{equation}
\Feas=\bigl\{\n\in\NN:\ \max(4,\lceil\nmin\rceil)\le\n\le\lfloor\nmax\rfloor
\bigr\},
\label{eq:window}
\end{equation}
and $\Feas\ne\emptyset$ if and only if $\max(4,\nmin)\le\nmax$. If in addition
$\Cost$ is nondecreasing, then $\nopt=\max(4,\lceil\nmin\rceil)$ and it is
unique.
\end{Theorem}

\begin{Proof}
By Lemma~\ref{lem:monotone}, $h$ is strictly decreasing, so
$\{\n:h(\n)\ge0\}$ is a down-set $(0,\nmax]$; combined with
Lemma~\ref{lem:detequiv} this shows constraint~(P) is equivalent to
$\n\le\nmax$. By Corollary~\ref{cor:nmin}, constraint~(S) is equivalent to
$\n\ge\nmin$. Intersecting the two half-lines with $\n\ge4$ and with $\NN$
gives~\eqref{eq:window}; a nonempty intersection of $[\nmin,\nmax]$ with
$[4,\infty)$ requires exactly $\max(4,\nmin)\le\nmax$. Minimizing a
nondecreasing $\Cost$ over a nonempty finite integer interval attains its
minimum at the left endpoint, and uniqueness holds when $\Cost$ is strictly
increasing.
\end{Proof}

\begin{Corollary}[Infeasibility certificate]
\label{cor:infeasible}
If $\nmin>\nmax$, no validator-set size satisfies the requirement. The system
must then be changed structurally --- by raising achievable concurrency~$\cc$,
by lengthening the detection window~$\w$, by relaxing $\epsp$ or $\epss$, or by
abandoning the flat BFT topology --- rather than by resizing the validator set.
\end{Corollary}

Corollary~\ref{cor:infeasible} is the practically useful half of
Theorem~\ref{thm:window}. Capacity-planning studies normally return a number;
here the procedure is also allowed to return the answer that no number exists,
together with the two inequalities that conflict. We are not aware of an
equivalent statement for permissioned BFT, though the underlying argument ---
intersecting two monotone constraints --- is elementary.

\begin{Remark}[Interpretation of the window]
The width $\nmax-\nmin$ measures the design margin. A narrow window means the
deployment is simultaneously close to its throughput ceiling and to its
detection floor, so small errors in the calibrated $\beta$ or in the measured
$\pd$ translate directly into infeasibility. Sensitivity of the window to
$\beta$, $\rho$ and network latency is reported in Section~\ref{sec:results}.
\end{Remark}

\subsection{Computation}

Both endpoints are available in closed form when $s$ is constant. In general
$\nmax$ is found by bisection, whose correctness is certified by
Lemma~\ref{lem:monotone}.

\begin{Algorithm}[!t]
\caption{Validator-set sizing}
\label{alg:sizing}
\REQUIRE Calibrated $(\Tzero,\gamma,\beta,\csat,\theta,\widehat{\Sigma},
         \hat\sigma)$; detector $(\pd,\pf,\rho)$; targets
         $\Dpk,\cc,\epsp,\epss,\w,\thr$; search bound $N$
\ENSURE $\nopt$, or an infeasibility certificate
\STATE $\varepsilon_{0}\leftarrow\mu_{B}-\thr$ \COMMENT{from~\eqref{eq:muB}}
\IF{$\varepsilon_{0}\le0$}
  \STATE \RET infeasible \COMMENT{detector cannot separate the attack}
\ENDIF
\STATE $\nmin\leftarrow\bigl(1+(\kf-1)\rho\bigr)\ln(1/\epss)\,/\,
        \bigl(2\w\varepsilon_{0}^{2}\bigr)$
\IF{$h(4)<0$}
  \STATE \RET infeasible \COMMENT{demand unmet even at the minimum set}
\ENDIF
\STATE $lo\leftarrow4$; $hi\leftarrow N$
\WHILE{$hi-lo>1$}
  \STATE $mid\leftarrow\lfloor(lo+hi)/2\rfloor$
  \IF{$h(mid)\ge0$}
    \STATE $lo\leftarrow mid$
  \ELSE
    \STATE $hi\leftarrow mid$
  \ENDIF
\ENDWHILE
\STATE $\nmax\leftarrow lo$
\IF{$\max(4,\lceil\nmin\rceil)>\nmax$}
  \STATE \RET infeasible \COMMENT{empty window, Corollary~\ref{cor:infeasible}}
\ENDIF
\STATE \RET $\nopt\leftarrow\max(4,\lceil\nmin\rceil)$
\end{Algorithm}

Algorithm~\ref{alg:sizing} performs $O(\log N)$ evaluations of~$h$, each of
which costs $O(1)$ arithmetic operations.
